\section{Availability, Reproducibility, and Limits}

The artifact repository contains the \tool Python package, adapter scripts,
metric runners, generated CSV/Markdown/\LaTeX{} tables, provenance logs,
BibTeX/source-audit notes, an MIT license, and a reviewer quickstart. The
reviewer entry point is the tag
\texttt{audit-sid-cikm-resource-v0.1}, with a clean-checkout verifier for the
published tables and figure. Large local artifacts are kept out of git and
described through manifests so reviewers can separate public reproducibility
assets from local caches.
The artifact ships with a clean-checkout verifier script that checks table
files, figure inputs, and claim-facing CSVs from manifest-pinned assets;
reviewer entry points are documented in the repository quickstart.

The resource is intended to be used in three modes. A paper reviewer can run
the clean-check verifier and compare published CSVs with the PDF tables. A
method author can add an adapter that emits the normalized
\texttt{sid\_assignments} table, then receive D1--D5a reports without changing
their training loop. A downstream recommender researcher can use the same
diagnostics as a pre-training gate: artifacts with missing items, inconsistent
depth, severe collisions, or weak collaborative-prefix recovery can be flagged
before generator training consumes GPU time. These modes are deliberately
lighter than a full benchmark suite, but they match the resource-track goal of
making a reusable software artifact available at review time.

The evidence supports a conservative resource-demo claim. The toolkit does not
reproduce TIGER end-to-end; the GRID Musical row is a controlled feature-text
export. The ReSID Sports balanced GAOQ path was stopped after constrained
k-means became CPU-bound, so the main ReSID evidence is the bounded Musical
export. CARD remains a fallback/stressor path, not faithful CARD
reproduction~\cite{wei2026card}. D2 includes an interaction-qualified
controller but not causal harm; D3 is not proven monotonic with Recall@K or
NDCG; D5a is structural rather than measured latency; and D6 is optional drift
evidence.

\tool audits tokenizer artifacts; it is not a new \sid tokenizer, a complete
generative-recommender evaluation, or an industrial online-impact study.
Artifact diagnostics can reveal pressure points before generator training, but
they do not replace ranking evaluation. Future versions should add causal
collision-harm validation, generator-output diagnostics, stronger same-dataset
method coverage, and unified search/recommendation checks~\cite{li2024survey,penha2025jointsid}.
Industrial SID deployment and ranking studies further motivate treating these
artifact properties as engineering objects~\cite{singh2023bettersemanticids,li2026sidcoord,ju2026snapchatsid}.

New adapters emit the same normalized tables, declare their evidence role, and
record feature or checkpoint provenance. Versioned manifests and fixed-name
``latest'' reports preserve stable entry points and timestamped provenance.
Thus the Musical case study is not a ranking of GRID against ReSID; it shows
how one artifact contract turns collisions, collaborative prefix recovery, tail
allocation, and prefix fan-out into separately checkable review objects.

The extension path is also explicit. A new named tokenizer improves method
coverage only when it exports item-level codes that pass the validator and can
be joined to metadata or interactions. A controlled stressor improves metric
calibration but does not count as method coverage. A generator trace activates
D7, but without such traces \tool reports only mapping-level and diagnostic
claims. This separation is what keeps the artifact useful as the SID literature
changes: new methods can be added without rewriting the diagnostics, while
claims remain tied to the evidence role of each row.
